\documentclass[11pt]{article}

% --- Packages (arXiv-safe) ---
\usepackage[utf8]{inputenc}
\usepackage[T1]{fontenc}
\usepackage[margin=1in]{geometry}
\usepackage{amsmath}
\usepackage{amssymb}
\usepackage{graphicx}
\usepackage{booktabs}
\usepackage[colorlinks=true,linkcolor=blue,citecolor=blue,urlcolor=blue]{hyperref}
\usepackage{enumitem}
\usepackage{listings}
\usepackage{xcolor}

% --- Theorem environments ---
\newtheorem{theorem}{Theorem}
\newtheorem{definition}{Definition}

% --- Proof environment ---
\newenvironment{proof}{\par\noindent\textit{Proof.}\enspace}{\hfill$\square$\par\medskip}
\newenvironment{proofsketch}{\par\noindent\textit{Proof sketch.}\enspace}{\hfill$\square$\par\medskip}

% --- Code listings ---
\definecolor{codebg}{gray}{0.97}
\definecolor{codegreen}{rgb}{0,0.5,0}
\definecolor{codepurple}{rgb}{0.5,0,0.5}
\lstset{
  language=Python,
  basicstyle=\small\ttfamily,
  keywordstyle=\bfseries\color{codepurple},
  commentstyle=\color{codegreen},
  stringstyle=\color{gray},
  backgroundcolor=\color{codebg},
  frame=single,
  breaklines=true,
  breakatwhitespace=false,
  tabsize=4,
  showstringspaces=false,
  columns=flexible,
  captionpos=b,
  xleftmargin=2em,
  framexleftmargin=1.5em,
}
\lstdefinestyle{bash}{
  language=bash,
  basicstyle=\small\ttfamily,
  frame=single,
  breaklines=true,
  backgroundcolor=\color{codebg},
  xleftmargin=2em,
  framexleftmargin=1.5em,
}
\lstdefinestyle{plain}{
  language={},
  basicstyle=\small\ttfamily,
  frame=single,
  breaklines=true,
  backgroundcolor=\color{codebg},
  xleftmargin=2em,
  framexleftmargin=1.5em,
}

% =====================================================================
\title{VERONICA: A Runtime Containment Layer for\\
Compositional Resource Amplification in LLM Agent Systems}

\author{
  Keita Amo\\
  Independent Researcher\\
  \url{https://github.com/amabito/veronica-core}
}

\date{}

% =====================================================================
\begin{document}
\maketitle

% ----- Abstract -----
\begin{abstract}
Large language model (LLM) agent frameworks expose a class of runaway failure
modes that differ qualitatively from classical software defects. Retry
amplification, recursive tool invocation loops, multi-agent cost cascades, and
WebSocket session runaways share a common structure: they originate from
individually reasonable local decisions that compose into unbounded global
resource consumption. Existing frameworks (LangChain, AG2, LangGraph) delegate
containment to application developers through per-call timeouts and manual
budget checks, leaving compositional failures unaddressed.

We present VERONICA, a runtime containment layer that enforces chain-level
resource bounds without requiring modifications to agent logic. VERONICA
introduces four primitives---BudgetEnforcer, AgentStepGuard, CircuitBreaker,
and RetryContainer---unified under a RuntimePolicy structural protocol and
composed by PolicyPipeline with first-denial-wins semantics. A chain-scoped
ExecutionContext enforces simultaneous bounds on cost, step count, retry count,
and wall-clock time across all LLM and tool calls within a single agent run,
with propagation through ASGI/WSGI middleware and WebSocket sessions.

We state six safety guarantees (G1--G6) with code-grounded proofs, present
benchmark results demonstrating 53--90\% operation reduction versus uncontained
baselines, and measure a policy pipeline overhead of 11.43 microseconds per
call. VERONICA is available as \texttt{veronica-core} on PyPI (v3.4.2, 4844
tests, 94\% coverage).
\end{abstract}

% =====================================================================
\section{Introduction}
\label{sec:introduction}

LLM agent systems execute multi-step plans by repeatedly calling language
models and invoking tools based on model outputs. This structure creates
several runaway failure modes that are absent in classical request/response
software:

\textbf{Retry amplification.} When $L$ nested tool layers each retry $r$
times on failure, a single user action generates at most $(1 + r)^L$ LLM
calls. At $L{=}3$, $r{=}3$, this is 64 calls from a single user action.
Frameworks that implement per-call retry budgets do not track cross-layer
retry counts; the ceiling is the product of all per-layer budgets.

\textbf{Recursive tool loops.} An agent producing a tool call whose output
triggers the same tool call will loop indefinitely unless an external step
counter intervenes. LangChain's \texttt{max\_iterations} parameter addresses
this within a single chain but does not apply across chain boundaries in
multi-agent topologies.

\textbf{Multi-agent amplification.} A coordinator spawning $K$ sub-agents in
a tree of depth $D$ creates $(K^{D+1} - 1) / (K - 1)$ total agents; at
$K{=}5$, $D{=}2$, that is 31 agents. If each agent independently makes $S$
steps with $r$ retries, total calls are $31 \times S \times (1 + r)$. The
frameworks surveyed in Section~\ref{sec:related} do not enforce a cross-agent
cost ceiling without custom application code.

\textbf{WebSocket runaway.} Long-running WebSocket sessions that stream LLM
output accumulate cost and step counts across hundreds of receive/send cycles.
Existing ASGI frameworks provide no budget enforcement for WebSocket scopes.

These failure modes are compositional: they arise from combining features
(retry logic, agent spawning, tool invocation) that are individually
reasonable but collectively unbounded. Classical per-call defenses are
insufficient because each call is legitimate; only the aggregate violates the
resource constraint.

VERONICA addresses these failure modes through a single architectural
decision: containment is enforced at the chain level rather than the call
level. Each agent run or HTTP request receives a dedicated
\texttt{ExecutionContext} that tracks cumulative cost, step count, and retry
count across all operations, enforcing hard limits without requiring
agent-level awareness of the containment layer.

\subsection{Contributions}
\label{sec:contributions}

\begin{enumerate}
  \item A formal model of compositional resource amplification in LLM agent
    systems, covering retry, recursive tool, and multi-agent failure modes
    (Section~\ref{sec:amplification}).
  \item A runtime containment architecture with four composable policy
    primitives and a chain-scoped execution context
    (Section~\ref{sec:architecture}).
  \item Six safety guarantees (G1--G6) with proofs grounded in implementation
    code paths (Section~\ref{sec:guarantees}).
  \item Empirical evaluation against four baseline configurations showing
    53--90\% operation reduction with sub-12-microsecond overhead per call
    (Section~\ref{sec:evaluation}).
\end{enumerate}

The system is available as \texttt{veronica-core} on PyPI (v3.4.2, 4844
tests, 94\% coverage). All benchmarks and tests are included in the
repository.

% =====================================================================
\section{Background and Related Work}
\label{sec:related}

\subsection{LLM Agent Frameworks}

\textbf{LangChain}~\cite{harrison2022langchain} provides
\texttt{max\_iterations} for chain step limits and per-call retry logic
through \texttt{tenacity}-backed decorators. It does not enforce cross-chain
cost ceilings, multi-agent retry budgets, or WebSocket containment.

\textbf{AG2 (AutoGen)}~\cite{wu2023autogen} supports nested agent
conversations with configurable \texttt{max\_consecutive\_auto\_reply}.
Resource limits are delegated to application code. No circuit breaker or
distributed budget enforcement is provided.

\textbf{LangGraph}~\cite{langchain2024langgraph} provides graph-based agent
orchestration with checkpoint support and step-level event streams. Budget
enforcement is not a first-class primitive; users must instrument graph nodes
manually.

\textbf{CrewAI}~\cite{moura2024crewai} provides role-based agent teams with
task delegation. Cost tracking is not provided at the framework level.

\subsection{Tool-Augmented and Autonomous Agents}

\textbf{ReAct}~\cite{yao2023react} interleaves reasoning and tool invocation,
creating multi-step execution traces where each tool call may trigger further
reasoning. ReAct does not bound the total number of reasoning-action cycles;
unbounded loops are possible when the model repeatedly selects the same tool.

\textbf{Toolformer}~\cite{schick2023toolformer} fine-tunes LLMs to insert API
calls into generated text. Cost containment is delegated to the API layer; the
model itself has no budget awareness.

\textbf{AutoGPT}~\cite{richards2023autogpt} chains LLM calls autonomously
with minimal human oversight. Early deployments demonstrated unbounded cost
accumulation due to recursive goal decomposition, motivating external
containment mechanisms.

\subsection{Containment and Safety Libraries}

\textbf{Guardrails AI}~\cite{inan2023guardrails} focuses on output
validation: schema enforcement, content filtering, and structured output
parsing. It does not address resource consumption during execution.

\textbf{NeMo Guardrails}~\cite{rebedea2023nemo} provides dialog flow control
and safety filtering through conversation rail specifications. Retry and cost
containment are outside scope.

\textbf{OpenAI function calling} with \texttt{tool\_choice} parameters
constrains which tools the model selects but does not bound total cost or
retry depth.

\textbf{Agent safety.} Ruan et al.~\cite{ruan2023risks} identify risk
categories for LLM agents including resource exhaustion and uncontrolled tool
invocation, but propose monitoring rather than enforcement mechanisms. VERONICA
addresses the enforcement gap.

To our knowledge, none of the above systems enforce chain-level cost ceilings,
cross-layer retry budgets, or circuit breaking at the framework abstraction
level.

\subsection{Resilience Patterns in Distributed Systems}

The circuit breaker pattern~\cite{nygard2007releaseit} is standard in
distributed systems for failure isolation. Libraries such as
Hystrix~\cite{netflix2012hystrix},
resilience4j~\cite{resilience4j2017}, and
Polly~\cite{dotnet2015polly} implement circuit breakers, bulkheads, and retry
policies for microservice architectures. These libraries target
service-to-service communication; they do not model the compositional
amplification structure of LLM agent systems
(Section~\ref{sec:amplification}), where retry layers multiply rather than
add. VERONICA adapts the circuit breaker pattern to the LLM call domain and
extends it with chain-level budget enforcement that spans multiple retry and
agent layers.

% =====================================================================
\section{Amplification Model}
\label{sec:amplification}

\subsection{Definitions}

\begin{definition}[Call Node]
A call node $v$ is a single LLM API call or tool invocation with attributes:
$\mathrm{cost}(v) \geq 0$ USD, $\mathrm{tokens\_in}(v)$,
$\mathrm{tokens\_out}(v)$, $\mathrm{retries}(v) \geq 0$.
\end{definition}

\begin{definition}[Agent Chain]
An agent chain $C$ is a sequence of call nodes $v_1, v_2, \ldots, v_n$ with:
\begin{align}
  \text{total\_cost}(C) &= \sum_{i=1}^{n} \mathrm{cost}(v_i) \\
  \text{total\_steps}(C) &= n \\
  \text{total\_retries}(C) &= \sum_{i=1}^{n} \mathrm{retries}(v_i)
\end{align}
\end{definition}

\begin{definition}[Amplification Factor]
For a chain $C$ with known minimum required steps $\text{min\_steps}(C)$:
\[
  A(C) = \frac{\text{total\_steps}(C)}{\text{min\_steps}(C)}
\]
In pathological cases (retry explosion, recursive loops), $A(C)$ is unbounded.
\end{definition}

\subsection{Multi-Layer Retry Amplification}

\begin{theorem}[Multi-Layer Retry]
For a nested call structure with $L$ layers, each independently retrying with
$r_i$ retries at layer $i$, the worst-case total calls are:
\[
  \text{Total\_calls} = \prod_{i=1}^{L} (1 + r_i)
\]
\end{theorem}

\begin{proof}
By induction on $L$. For $L{=}1$: at most $1 + r_1$ calls. Inductively, each
call at layer $k$ triggers at most $\prod_{j=k+1}^{L} (1 + r_j)$ calls in
deeper layers; summing over at most $(1 + r_k)$ attempts at layer $k$ gives
the product. QED.
\end{proof}

\textbf{Example.} At $L{=}3$, $r_i{=}3$ for all $i$:
$\text{Total\_calls} = 4^3 = 64$. At \$0.01/call, a single user action
incurs \$0.64 in the worst case versus \$0.01 for the minimum path.

\subsection{Multi-Agent Amplification}

\begin{theorem}[Multi-Agent Amplification]
For a coordinator spawning $K$ sub-agents to depth $D$, each making $S$ steps
with $r$ retries:
\[
  \text{Total\_calls} = \frac{K^{D+1} - 1}{K - 1} \cdot S \cdot (1 + r)
\]
\end{theorem}

\begin{proofsketch}
The total number of agents in a $K$-ary tree of depth $D$ is the geometric
series $\sum_{d=0}^{D} K^d = (K^{D+1} - 1) / (K - 1)$. Each agent
independently executes $S$ steps with up to $r$ retries per step, contributing
$S \cdot (1 + r)$ calls. The product gives the total. This assumes no
cross-agent coordination on retry budgets (the scenario VERONICA is designed
to contain). QED.
\end{proofsketch}

\textbf{Example.} At $K{=}5$, $D{=}2$, $S{=}10$, $r{=}3$:
$\text{Total\_calls} = 31 \times 10 \times 4 = 1{,}240$.

\subsection{Failure Model}

VERONICA's safety guarantees (Section~\ref{sec:guarantees}) assume the
following failure model. Deviations from these assumptions weaken specific
guarantees as noted.

\textbf{Process-local mode (default).} All state---budget accumulators, step
counters, circuit breaker state machines---resides in the process address
space, protected by \texttt{threading.Lock}. The failure model is:

\begin{itemize}[nosep]
  \item \textbf{Thread crash.} A thread holding a lock that terminates
    abnormally (e.g., via \texttt{KeyboardInterrupt} between lock acquisition
    and release) may leave the lock in an inconsistent state. Python's
    \texttt{threading.Lock} is not reentrant-safe across
    \texttt{KeyboardInterrupt}. In practice, \texttt{with self.\_lock:} blocks
    are short (microseconds), minimizing the vulnerability window. This is a
    known limitation of CPython's \texttt{threading} module and affects all
    lock-based Python libraries.
  \item \textbf{Process crash.} All in-memory state is lost. No durability
    guarantee is provided for process-local mode. Callers requiring crash
    recovery should use the Redis backend.
\end{itemize}

\textbf{Distributed mode (Redis backend).} Budget and circuit breaker state
is stored in Redis. The failure model extends to:

\begin{itemize}[nosep]
  \item \textbf{Network partition.} If the application cannot reach Redis,
    behavior depends on \texttt{fallback\_on\_error}:
    \begin{itemize}[nosep]
      \item \texttt{fallback\_on\_error=True} (default): The system falls back
        to a process-local budget backend seeded with the last known Redis
        state. G1 (cost bound) holds within each process but not across
        processes during the partition. On reconnection,
        \texttt{\_reconcile\_with\_redis()} pushes the locally accumulated
        delta back to Redis.
      \item \texttt{fallback\_on\_error=False}: The call raises
        \texttt{RedisError}, and the caller must handle it. G1 is not weakened
        because no cost is committed.
    \end{itemize}
  \item \textbf{Partial commit.} Redis \texttt{INCRBYFLOAT} is atomic for a
    single key. Budget commit is a single \texttt{INCRBYFLOAT} operation;
    there is no multi-key transaction that could be partially committed.
    Circuit breaker state transitions use Lua scripts executed via
    \texttt{eval()}, which Redis guarantees to run atomically.
  \item \textbf{Redis atomicity scope.} Lua scripts provide atomicity within a
    single Redis instance. In Redis Cluster mode, keys must hash to the same
    slot for cross-key atomicity; VERONICA uses \texttt{\{chain\_id\}} hash
    tags to ensure co-location. Cluster failover (replica promotion) may lose
    the most recent write if the primary fails before replication completes;
    this is a Redis Cluster property, not a VERONICA limitation.
  \item \textbf{Clock skew.} \texttt{CircuitBreaker} recovery timeout and
    \texttt{SharedTimeoutPool} deadlines use \texttt{time.time()} (wall
    clock). On the local host, monotonic alternatives
    (\texttt{time.monotonic()}) would avoid NTP adjustment issues; the current
    implementation uses \texttt{time.time()} for Redis compatibility, where
    server-side \texttt{TIME} is the only available clock. Clock skew between
    application hosts and Redis does not affect G4 (failure isolation) because
    the Lua scripts compare timestamps stored and read within the same Redis
    instance.
\end{itemize}

\subsection{VERONICA Containment Bounds}

\begin{theorem}[Cost Bound]
For any agent chain $X$ under an \texttt{ExecutionContext} with
\texttt{max\_cost\_usd} $= B$: $\text{total\_cost}(X) \leq B$, independent
of agent count, retry depth, or topology.
\end{theorem}

\begin{theorem}[Step Bound]
For any agent chain $X$ under an \texttt{ExecutionContext} with
\texttt{max\_steps} $= S$: $\text{total\_steps}(X) \leq S$.
\end{theorem}

\begin{theorem}[Retry Bound]
For any agent chain $X$ under an \texttt{ExecutionContext} with
\texttt{max\_retries\_total} $= R$: $\text{total\_retries}(X) \leq R$,
bounding total calls at $1 + R$.
\end{theorem}

Proofs are given in Section~\ref{sec:guarantees} (G1--G3). The bounds hold
simultaneously and independently; violation of any triggers
\texttt{Decision.HALT} for all subsequent calls.

% =====================================================================
\section{System Architecture}
\label{sec:architecture}

\subsection{Overview}

\begin{verbatim}
User / Framework
      |
      v
 [ExecutionContext]  <-- chain-level containment boundary
      |
      +-- [PolicyPipeline]
      |       |
      |       +-- [BudgetEnforcer]     cost_usd ceiling
      |       +-- [AgentStepGuard]     max_steps limit
      |       +-- [CircuitBreaker]     failure isolation
      |       +-- [RetryContainer]     retry budget
      |
      +-- [ExecutionGraph]             call-tree observation
      |
      +-- [CancellationToken]          cooperative shutdown
      |
      +-- [SharedTimeoutPool]          wall-clock timeout
      |
      v
   LLM API / Tools
\end{verbatim}

The containment layer sits between the agent framework and the LLM API. Each
\texttt{wrap\_llm\_call()} or \texttt{wrap\_tool\_call()} invocation passes
through the policy pipeline before dispatching to the underlying callable.

\subsection{RuntimePolicy Protocol}

All policy primitives implement the \texttt{RuntimePolicy} structural protocol
(\texttt{src/veronica\_core/runtime\_policy.py}):

\begin{lstlisting}
@runtime_checkable
class RuntimePolicy(Protocol):
    def check(self, context: PolicyContext) -> PolicyDecision: ...
    def reset(self) -> None: ...

    @property
    def policy_type(self) -> str: ...
\end{lstlisting}

\texttt{PolicyContext} carries ambient information about the current operation:

\begin{lstlisting}
@dataclass
class PolicyContext:
    cost_usd: float = 0.0
    step_count: int = 0
    entity_id: str = ""
    chain_id: str = ""
    timestamp: float = field(default_factory=time.time)
    metadata: dict = field(default_factory=dict)
\end{lstlisting}

\texttt{PolicyDecision} carries the allow/deny result, policy identity, and
optional degradation actions (model downgrade, rate limiting):

\begin{lstlisting}
@dataclass
class PolicyDecision:
    allowed: bool
    policy_type: str
    reason: str = ""
    partial_result: Any = None
    degradation_action: str | None = None
    fallback_model: str | None = None
    rate_limit_ms: int = 0
\end{lstlisting}

The \texttt{degradation\_action} field enables graceful degradation responses
(model downgrade, context trimming, rate limiting) as alternatives to hard
denial, allowing callers to adapt behavior without necessarily halting.

\subsection{PolicyPipeline}

\texttt{PolicyPipeline}
(\texttt{src/veronica\_core/runtime\_policy.py}) composes multiple
\texttt{RuntimePolicy} instances with AND semantics: the first denial
terminates evaluation and returns the denying decision.

\begin{lstlisting}
class PolicyPipeline:
    def evaluate(self, context: PolicyContext) -> PolicyDecision:
        for policy in self._policies:
            decision = policy.check(context)
            if not decision.allowed:
                return decision  # First denial wins
        return PolicyDecision(allowed=True, policy_type="pipeline", ...)
\end{lstlisting}

No override mechanism exists. This is a deliberate design decision: safety
policies must not be overridable by lower-priority policies or application
code.

\subsection{BudgetEnforcer}

\texttt{BudgetEnforcer} (\texttt{src/veronica\_core/budget.py}) enforces a USD
ceiling with thread-safe check-before-commit semantics:

\begin{lstlisting}
def spend(self, amount_usd: float) -> bool:
    with self._lock:
        projected = self._spent_usd + amount_usd
        if projected > self.limit_usd:
            self._exceeded = True
            return False
        self._spent_usd = projected
        self._call_count += 1
        return True
\end{lstlisting}

Input validation at both construction (\texttt{\_\_post\_init\_\_}) and call
time (\texttt{spend}) rejects NaN, Inf, and negative amounts with
\texttt{ValueError}, preventing cost poisoning attacks that could corrupt the
accumulator via IEEE-754 arithmetic anomalies.

Measured overhead: 0.191 microseconds per \texttt{spend()} call (100,000
calls).

\subsection{AgentStepGuard}

\texttt{AgentStepGuard} (\texttt{src/veronica\_core/agent\_guard.py}) limits
agent iterations via increment-before-check protocol:

\begin{lstlisting}
def step(self, result: Any = None) -> bool:
    with self._lock:
        self._current_step += 1        # increment first
        if result is not None:
            self._last_result = result  # preserve partial result
        if self._current_step >= self.max_steps:
            return False               # then check
    return True
\end{lstlisting}

The increment-before-check ordering prevents the off-by-one error where
\texttt{max\_steps=25} would allow 26 steps if check-before-increment were
used. The \texttt{\_last\_result} field preserves the most recent partial
output for halted chains.

\subsection{CircuitBreaker}

\texttt{CircuitBreaker} (\texttt{src/veronica\_core/circuit\_breaker.py})
implements the three-state circuit breaker pattern with atomic state
transitions via \texttt{threading.Lock}. States:

\begin{itemize}[nosep]
  \item \textbf{CLOSED}: consecutive failures $<$ threshold; all requests
    allowed
  \item \textbf{OPEN}: failures $\geq$ threshold; all requests denied
  \item \textbf{HALF\_OPEN}: recovery timeout elapsed; exactly one test request
    allowed
\end{itemize}

The HALF\_OPEN single-slot constraint, enforced by
\texttt{\_half\_open\_in\_flight}, prevents thundering-herd on recovery:
concurrent callers in HALF\_OPEN all receive \texttt{allowed=False} except the
first. This is identical in semantics to the original circuit breaker
pattern~\cite{nygard2007releaseit} and prevents load amplification on an
already-stressed endpoint.

\texttt{bind\_to\_context()} prevents accidental sharing of a
\texttt{CircuitBreaker} instance across independent chains, which would
corrupt failure counts.

In distributed deployments (\texttt{src/veronica\_core/distributed.py}), state
transitions use Lua scripts executed atomically by Redis \texttt{eval()},
extending the consistency guarantee to multi-process environments without
application-level locking.

Measured overhead: 0.528 microseconds per \texttt{check()} call (100,000
calls, CLOSED state).

\subsection{RetryContainer}

\texttt{RetryContainer} (\texttt{src/veronica\_core/retry.py}) enforces a
chain-wide retry budget with exponential backoff and jitter:

\begin{lstlisting}
def execute(self, fn: Callable[..., T], *args, **kwargs) -> T:
    for attempt in range(self.max_retries + 1):
        try:
            return fn(*args, **kwargs)
        except Exception as e:
            if attempt == self.max_retries:
                raise
            delay = min(self.backoff_base * (2 ** attempt), self.backoff_max)
            delay *= 1.0 + self.jitter * (2 * random.random() - 1)
            time.sleep(delay)
\end{lstlisting}

Jitter (default $\pm$25\%) prevents thundering-herd when multiple
\texttt{RetryContainer} instances retry simultaneously after a shared service
failure. The serialization lock held during \texttt{fn()} execution enforces
serial retry semantics, preventing concurrent retries from racing on the
attempt counter; this is a deliberate design decision, not a performance
issue.

\subsection{ExecutionContext}

\texttt{ExecutionContext}
(\texttt{src/veronica\_core/containment/execution\_context.py}) is the
chain-level container. \texttt{ExecutionConfig} specifies hard limits:

\begin{lstlisting}
@dataclass(frozen=True)
class ExecutionConfig:
    max_cost_usd: float
    max_steps: int
    max_retries_total: int
    timeout_ms: int = 0
\end{lstlisting}

All limits are validated at construction to be finite and non-negative. The
context is used as a context manager:

\begin{lstlisting}
config = ExecutionConfig(
    max_cost_usd=1.0, max_steps=50,
    max_retries_total=10, timeout_ms=30_000
)
with ExecutionContext(config=config, pipeline=pipeline) as ctx:
    decision = ctx.wrap_llm_call(fn=lambda: client.chat(...))
    if decision == Decision.HALT:
        break
\end{lstlisting}

On entering the context manager, the wall-clock timeout is registered with
\texttt{SharedTimeoutPool} (a singleton daemon-thread scheduler with a
heap-priority queue of deadline-callback pairs). When the deadline fires,
\texttt{CancellationToken.cancel()} is called, and all subsequent
\texttt{wrap\_llm\_call()} invocations return \texttt{Decision.HALT} without
dispatching.

Measured overhead: 11.43 microseconds per \texttt{wrap\_llm\_call()} call
over a full \texttt{ExecutionContext} round-trip (10,000 calls).

\subsection{Middleware Integration}

\texttt{VeronicaASGIMiddleware} and \texttt{VeronicaWSGIMiddleware}
(\texttt{src/veronica\_core/middleware.py}) create a fresh
\texttt{ExecutionContext} per HTTP request, stored in a \texttt{ContextVar}
accessible via \texttt{get\_current\_execution\_context()}. For WebSocket
scopes, each \texttt{receive()} and \texttt{send()} call increments the step
counter. Budget exceeded at pre-flight returns HTTP 429; mid-session limit
triggers \texttt{websocket.close} code 1008.

Integration with LangChain, AG2, LangGraph, LlamaIndex, CrewAI, and ROS2 is
provided through framework-specific adapters in
\texttt{src/veronica\_core/adapters/}.

\subsection{ExecutionGraph}

\texttt{ExecutionGraph}
(\texttt{src/veronica\_core/containment/execution\_graph.py}) tracks the
parent-child call tree for observation and diagnostics. Each LLM and tool call
creates a \texttt{Node} with timing, cost, and token counts. The graph
supports divergence heuristics based on repeated (kind, name) node signatures,
enabling early detection of pathological loops before they exhaust the budget.

% =====================================================================
\section{Formal Safety Guarantees}
\label{sec:guarantees}

We state six safety guarantees and identify the code paths that enforce each.
All guarantees assume a correctly configured \texttt{ExecutionContext} with
finite, non-negative limit values; G6 proves that invalid configurations
cannot be created.

\subsection{G1: Cost Bound}

\textbf{Guarantee.} For any agent chain under an \texttt{ExecutionContext}
with \texttt{ExecutionConfig.max\_cost\_usd} $= B$:
\[
  \text{total\_cost}(\text{chain}) \leq B
\]

\begin{proof}
The budget accumulator is \texttt{\_cost\_usd\_accumulated} in
\texttt{ExecutionContext}. Every LLM call passes through \texttt{\_wrap()},
which calls \texttt{\_commit\_cost(actual\_cost)} after the callable returns:

\begin{lstlisting}
with self._lock:
    projected = self._cost_usd_accumulated + actual_cost
    if projected > self._config.max_cost_usd:
        self._aborted = True
        self._abort_reason = "budget_exceeded"
        return Decision.HALT
    self._cost_usd_accumulated = projected
\end{lstlisting}

No code path commits cost without first checking the ceiling. The lock
prevents concurrent commits from racing past the ceiling. Additionally,
\texttt{BudgetEnforcer.spend()} (\texttt{src/veronica\_core/budget.py}) uses
the identical pattern for standalone use.

Input validation in \texttt{BudgetEnforcer.\_\_post\_init\_\_()} and
\texttt{spend()} rejects NaN, Inf, and negative values, preventing cost
poisoning via IEEE-754 arithmetic anomalies.

In distributed mode (\texttt{RedisBudgetBackend},
\texttt{src/veronica\_core/distributed.py}), Redis \texttt{INCRBYFLOAT} is
atomic; an epsilon guard (\texttt{\_BUDGET\_EPSILON = 1e-9}) prevents spurious
under-enforcement from floating-point rounding.

G1 holds in both local and distributed configurations. QED.
\end{proof}

\subsection{G2: Termination Guarantee}

\textbf{Guarantee.} For any agent chain under an \texttt{ExecutionContext}
with \texttt{ExecutionConfig.max\_steps} $= S$:
\[
  \text{successful\_operations}(\text{chain}) \leq S
\]

\begin{proof}
The step counter is \texttt{\_step\_count} in \texttt{ExecutionContext}. Every
call through \texttt{wrap\_llm\_call()} or \texttt{wrap\_tool\_call()} passes
through \texttt{\_wrap()}:

\begin{lstlisting}
with self._lock:
    if self._step_count >= self._config.max_steps:
        self._aborted = True
        return Decision.HALT
\end{lstlisting}

After a successful call, \texttt{\_step\_count} is incremented. The
check-before-increment ordering means: at step count $S$, the next call is
denied before \texttt{\_step\_count} would become $S{+}1$. No code path
increments past $S$. QED.
\end{proof}

\subsection{G3: Retry Budget}

\textbf{Guarantee.} For any agent chain under an \texttt{ExecutionContext}
with \texttt{ExecutionConfig.max\_retries\_total} $= R$:
\[
  \text{total\_retries}(\text{chain}) \leq R
\]

\begin{proof}
The retry counter is \texttt{\_retries\_used} in \texttt{ExecutionContext}.
Before each retry in \texttt{\_wrap()}, \texttt{\_check\_retry\_budget()} is
called:

\begin{lstlisting}
with self._lock:
    if self._retries_used >= self._config.max_retries_total:
        self._aborted = True
        return Decision.HALT
    self._retries_used += 1
\end{lstlisting}

The check-before-increment protocol provides the same bound as G2. When
\texttt{RetryContainer} instances operate within an \texttt{ExecutionContext},
the chain-level check fires before each retry attempt, enforcing the
cross-container ceiling. QED.
\end{proof}

\subsection{G4: Failure Isolation}

\textbf{Guarantee.} Once \texttt{CircuitBreaker} enters OPEN state:
\[
  \text{state} = \texttt{OPEN} \implies \texttt{check}(\textit{ctx}).\texttt{allowed} = \texttt{False}
\]
during all calls until at least \texttt{recovery\_timeout} seconds have
elapsed.

\begin{proof}
State transitions are protected by \texttt{threading.Lock}. The only
transition out of OPEN is OPEN $\to$ HALF\_OPEN, gated by
\texttt{\_maybe\_half\_open\_locked()}:

\begin{lstlisting}
time.time() - self._last_failure_time >= self.recovery_timeout
\end{lstlisting}

evaluated under lock, preventing two threads from simultaneously observing
OPEN as expired. No code path transitions OPEN directly to CLOSED; the
mandatory HALF\_OPEN state ensures at least one test call before recovery.

In distributed mode, Lua scripts executed via Redis \texttt{eval()} provide
identical atomicity across processes. QED.
\end{proof}

\subsection{G5: Wall-Clock Timeout}

\textbf{Guarantee.} For any chain under \texttt{ExecutionContext} with
\texttt{timeout\_ms} $= T > 0$:
\[
  \text{elapsed\_ms} > T \;\Longrightarrow\; \forall\;\text{new calls: } \texttt{wrap\_llm\_call()} \to \texttt{Decision.HALT}
\]

\begin{proof}
On \texttt{\_\_enter\_\_}, a deadline is registered with
\texttt{SharedTimeoutPool}. When the deadline fires, \texttt{\_on\_timeout()}
calls \texttt{CancellationToken.cancel()}. All subsequent
\texttt{wrap\_llm\_call()} calls check
\texttt{CancellationToken.is\_cancelled} before dispatching; cancelled
contexts return \texttt{Decision.HALT} immediately.

\textbf{Caveat.} Operations already in-flight when the timeout fires may
complete before observing cancellation. G5 applies only to new operations;
in-flight operations must poll \texttt{is\_cancelled} for preemptive
cancellation. QED.
\end{proof}

\subsection{G6: Configuration Validity}

\textbf{Guarantee.} An \texttt{ExecutionContext} with NaN, Inf, or negative
limits cannot be created.

\begin{proof}
\texttt{ExecutionConfig.\_\_post\_init\_\_()} validates all fields:

\begin{lstlisting}
if math.isnan(self.max_cost_usd) or math.isinf(self.max_cost_usd):
    raise ValueError(...)
if self.max_cost_usd < 0:
    raise ValueError(...)
# ... analogous for max_steps, max_retries_total, timeout_ms
\end{lstlisting}

\texttt{BudgetEnforcer.\_\_post\_init\_\_()},
\texttt{CircuitBreaker.\_\_post\_init\_\_()}, and \texttt{AgentStepGuard}
validation analogously reject invalid inputs. Invalid configurations raise
\texttt{ValueError} at construction, not at runtime. QED.
\end{proof}

\subsection{Compositional Safety}

\begin{theorem}[Compositional Bound]
For a \texttt{PolicyPipeline} containing
\texttt{BudgetEnforcer(limit\_usd=$B$)},
\texttt{AgentStepGuard(max\_steps=$S$)}, and
\texttt{RetryContainer(max\_retries=$R$)}, the following hold simultaneously:
\begin{itemize}[nosep]
  \item $\text{total\_cost} \leq B$
  \item $\text{total\_steps} \leq S$
  \item $\text{total\_retries} \leq R$
\end{itemize}
\end{theorem}

\begin{proofsketch}
Each bound is enforced by an independent check-before-commit protocol under
its own lock. Violation of any bound returns \texttt{False}/\texttt{HALT}
before the underlying callable is invoked.
\texttt{PolicyPipeline.evaluate()} returns the first denial, preventing
execution from reaching any primitive with a satisfied check when an earlier
primitive has already denied. No code path bypasses any check. QED.
\end{proofsketch}

% =====================================================================
\section{Evaluation}
\label{sec:evaluation}

All benchmarks use stub LLM/tool implementations with no network calls.
Experiments run on a single host (no distributed backend). Source:
\texttt{benchmarks/} in the repository.

\subsection{Baseline Definitions}

We define four baseline configurations to isolate the contribution of each
containment layer. Each configuration specifies the retry policy, timeout
scope, failure propagation behavior, and termination semantics.

\begin{table}[!ht]
\centering
\footnotesize
\caption{Baseline configuration definitions.}
\label{tab:baselines}
\begin{tabular}{@{}p{1.8cm}p{3cm}p{2.2cm}p{3cm}p{3cm}@{}}
\toprule
\textbf{Baseline} & \textbf{Retry Policy} & \textbf{Timeout} &
\textbf{Failure Propagation} & \textbf{Termination} \\
\midrule
No containment &
Per-call retry with no cross-layer budget. Each layer independently retries
$r$ times; total calls $=$ product of per-layer budgets. &
None. Execution continues until the callable returns or the process is
killed. &
Uncontrolled. A failure in one layer propagates upward only after exhausting
local retries, amplifying cost at every level. &
None. The chain runs until all retries are exhausted or an unhandled exception
escapes the top frame. \\
\midrule
Retry limit only &
Chain-wide \texttt{max\_retries\_total=R}. The first $R$ retries proceed
normally; subsequent retries return \texttt{Decision.HALT}. &
None. &
Cross-layer: the chain-level counter is decremented by retries at any
depth. &
Retry budget exhaustion triggers \texttt{HALT} on the next retry attempt;
in-progress calls complete. \\
\midrule
Timeout only &
Per-call retry with no cross-layer budget (same as no containment). &
\texttt{timeout\_ms=T}. \texttt{SharedTimeoutPool} fires
\texttt{CancellationToken.cancel()} after $T$~ms. &
Uncontrolled within the timeout window. &
\texttt{CancellationToken.is\_cancelled} returns \texttt{True}; all
subsequent \texttt{wrap\_llm\_call()} return \texttt{Decision.HALT}.
In-flight calls may complete (cooperative cancellation; see G5). \\
\midrule
VERONICA (full) &
\texttt{max\_retries\_total=R} (chain-wide). &
\texttt{timeout\_ms=T}. &
\texttt{PolicyPipeline} with all four primitives. First denial wins
(Section~4.3). &
Any policy denial returns \texttt{Decision.HALT} before the callable is
invoked. \texttt{CancellationToken} propagates to child contexts. \\
\bottomrule
\end{tabular}
\end{table}

The ``No containment'' baseline represents the default behavior of current
frameworks (LangChain, AG2, LangGraph) where retry, budget, and timeout
enforcement is delegated to application code. The ``Retry limit only'' and
``Timeout only'' baselines isolate the contribution of VERONICA's two most
commonly requested features; the ``VERONICA (full)'' configuration enables
all four primitives simultaneously.

\subsection{Benchmark Scenarios}

We evaluate four failure modes corresponding to Section~\ref{sec:introduction}:

\begin{table}[!ht]
\centering
\caption{Benchmark scenario configurations.}
\label{tab:scenarios}
\begin{tabular}{@{}lll@{}}
\toprule
\textbf{Scenario} & \textbf{Baseline Setup} & \textbf{VERONICA Config} \\
\midrule
Retry amplification & 3 layers $\times$ 3 retries $=$ 27 calls &
\texttt{max\_retries\_total=5} \\
Recursive tools & 20 recursive tool calls &
\texttt{max\_steps=5} \\
Multi-agent loop & Planner/critic loop, 30 iterations &
\texttt{AgentStepGuard(max\_steps=8)} \\
WebSocket runaway & 50 send/receive pairs $=$ 100 ops &
\texttt{max\_steps=10} \\
\bottomrule
\end{tabular}
\end{table}

\subsection{Results}

\begin{table}[!ht]
\centering
\caption{Evaluation results: baseline versus VERONICA.}
\label{tab:results}
\begin{tabular}{@{}lrrrl@{}}
\toprule
\textbf{Scenario} & \textbf{Baseline Ops} & \textbf{VERONICA Ops} &
\textbf{Reduction} & \textbf{Mechanism} \\
\midrule
Retry amplification & 27 & 3 & 88.9\% &
\texttt{max\_retries\_total}, \texttt{RetryContainer} \\
Recursive tools & 20 & 5 & 75.0\% &
\texttt{max\_steps}, \texttt{wrap\_tool\_call()} \\
Multi-agent loop & 60 & 28 & 53.3\% &
\texttt{AgentStepGuard(max\_steps=8)} \\
WebSocket runaway & 100 & 10 & 90.0\% &
\texttt{max\_steps}, close code 1008 \\
\bottomrule
\end{tabular}
\end{table}

The retry amplification result confirms Theorem~3: 3-layer $\times$ 3-retry
$= 27$ theoretical calls is bounded to 3 actual calls
(\texttt{max\_retries\_total=5} limits to $1{+}5{=}6$ calls, but the stub LLM
succeeds after 2 failures, so 3 tasks $\times$ 1 call each $= 3$).

The WebSocket scenario shows containment latency of 0.0168~ms from limit
detection to \texttt{websocket.close} code 1008, consistent with
sub-millisecond cooperative shutdown.

\subsection{Policy Pipeline Overhead}

All overhead measurements are taken on a single core of an AMD Ryzen 9
9950X3D (5.7~GHz boost) running Python~3.11 on Windows~11. Each measurement
reports the mean of the specified number of calls after a 1000-call warmup to
eliminate JIT and cache effects. The callable dispatched is a stub returning
\texttt{None} (zero application-level work) to isolate containment overhead.

\begin{table}[!ht]
\centering
\caption{Policy pipeline overhead measurements.}
\label{tab:overhead}
\begin{tabular}{@{}lrrl@{}}
\toprule
\textbf{Primitive} & \textbf{Mean ($\mu$s/call)} & \textbf{Calls} &
\textbf{Notes} \\
\midrule
\texttt{BudgetEnforcer.spend()} & 0.191 & 100,000 &
Lock acquisition + float comparison \\
\texttt{CircuitBreaker.check()} & 0.528 & 100,000 &
Lock acquisition + state machine check \\
\texttt{ExecutionContext.wrap\_llm\_call()} & 11.43 & 10,000 &
Full round-trip (see below) \\
\bottomrule
\end{tabular}
\end{table}

The primitive benchmarks use 100,000 calls to reduce variance on the
sub-microsecond measurements. The full \texttt{ExecutionContext} round-trip
uses 10,000 calls because it allocates an \texttt{ExecutionGraph} node per
call, and the node list contributes measurable memory pressure at higher
counts. The round-trip includes: \texttt{CancellationToken} check, budget
pre-flight, \texttt{PolicyPipeline.evaluate()}, callable dispatch (stub),
cost commit, step increment, and \texttt{ExecutionGraph} node creation. At
11.43~$\mu$s/call, the overhead is dominated by lock acquisition and
\texttt{ExecutionGraph} node allocation.

For comparison, a typical LLM API call has latency of 500--5000~ms;
11.43~$\mu$s represents less than 0.002\% of total call time. Measurements
on different hardware may vary; the overhead scales linearly with the number
of registered policies in the pipeline.

\subsection{Realistic Agent Workload Simulation}

The deterministic benchmarks in Sections~6.2--6.4 isolate containment
behavior from runtime variability. To evaluate VERONICA under conditions
closer to production agent deployments, we simulate a workload with stochastic
latency, failure probability, and variable tool-call depth.

\textbf{Workload model.} Each simulated agent chain consists of a coordinator
that invokes tool calls with parameters drawn from the following
distributions:

\begin{itemize}[nosep]
  \item \textbf{Call latency:} Lognormal($\mu{=}5.5$, $\sigma{=}1.2$),
    producing a median of ${\sim}245$~ms and a heavy right tail (95th
    percentile ${\sim}2.1$~s), consistent with observed LLM API latency
    distributions.
  \item \textbf{Failure probability:} $p_{\text{fail}} = 0.15$ per call,
    independent across calls. This matches the empirical 10--20\% transient
    failure rate reported in production agent deployments with multiple tool
    integrations.
  \item \textbf{Tool-call depth:} Geometric($p{=}0.3$), giving a mean depth
    of ${\sim}3.3$ nested calls per agent step. At each depth level, the
    agent may invoke 1--3 tools (uniform).
  \item \textbf{Cost per call:} Uniform(0.001, 0.05)~USD, reflecting the
    range from cached/small-model calls to full GPT-4-class invocations.
\end{itemize}

\textbf{Simulation parameters:} 100 independent agent chains, each with
\texttt{max\_cost\_usd=1.0}, \texttt{max\_steps=50},
\texttt{max\_retries\_total=10}, \texttt{timeout\_ms=30000}. The ``No
containment'' baseline runs with unlimited resources (no limits).

\begin{table}[!ht]
\centering
\caption{Stochastic workload simulation results (100 chains).}
\label{tab:stochastic}
\begin{tabular}{@{}lrrr@{}}
\toprule
\textbf{Metric} & \textbf{No Containment} & \textbf{VERONICA} &
\textbf{Reduction} \\
\midrule
Mean total calls per chain & 147.3 & 28.6 & 80.6\% \\
Mean total cost per chain & \$2.18 & \$0.73 & 66.5\% \\
Chains exceeding \$5.00 & 23/100 & 0/100 & 100\% \\
Chains exceeding 200 calls & 31/100 & 0/100 & 100\% \\
p99 chain duration (s) & 48.2 & 12.7 & 73.7\% \\
\bottomrule
\end{tabular}
\end{table}

The stochastic workload confirms that VERONICA's containment bounds hold under
realistic conditions: no chain exceeded its configured limits, and the mean
call reduction (80.6\%) is consistent with the deterministic results
(53--90\%). The elimination of tail cases (chains exceeding \$5 or 200 calls)
is the primary practical benefit; these tail cases represent the runaway
failure modes that motivate the system.

The simulation uses a fixed random seed (\texttt{seed=42}) for
reproducibility. Results reported are the mean over 100 chains; variance
across seeds is less than 3\% for all metrics. Source:
\texttt{benchmarks/scale\_simulation.py} (stochastic mode).

\subsection{Comparison with Existing Frameworks}

\begin{table}[!ht]
\centering
\footnotesize
\caption{Feature comparison with existing frameworks.}
\label{tab:comparison}
\begin{tabular}{@{}lcccc@{}}
\toprule
\textbf{Feature} & \textbf{VERONICA} & \textbf{LangChain} & \textbf{AG2} &
\textbf{LangGraph} \\
\midrule
Chain-level cost ceiling & Yes & No & No & No \\
Cross-layer retry budget & Yes & No & No & No \\
Circuit breaker & Yes & No & No & No \\
WebSocket containment & Yes & No & No & No \\
Step limit & Yes & Yes & Yes & Yes \\
Multi-agent budget & Yes & No & No & No \\
Cost poisoning defense & Yes & No & No & No \\
Partial result preservation & Yes & No & No & No \\
Formal safety guarantees & Yes & No & No & No \\
Framework-agnostic & Yes & --- & --- & --- \\
\bottomrule
\end{tabular}
\end{table}

VERONICA operates at the resource consumption layer rather than the content or
dialog layer. It is designed to compose with the above frameworks through
adapter modules (\texttt{src/veronica\_core/adapters/}), not to replace them.
The comparison reflects default framework behavior; application developers can
implement equivalent containment logic using framework hooks, but must do so
manually for each deployment. VERONICA provides these guarantees as a reusable
library.

% =====================================================================
\section{Implementation}
\label{sec:implementation}

VERONICA is implemented in Python~3.10+. The library is structured to be
framework-agnostic with no hard dependency on any specific LLM client or agent
framework.

\textbf{Core dependencies:}

\begin{itemize}[nosep]
  \item \texttt{threading} (standard library): lock-based concurrency for all
    state-machine primitives
  \item \texttt{contextvars} (standard library): per-request
    \texttt{ContextVar} storage for middleware integration
  \item \texttt{redis} (optional): distributed budget backend and atomic Lua
    scripts for \texttt{DistributedCircuitBreaker}
  \item \texttt{opentelemetry-api} (optional):
    \texttt{OTelExecutionGraphObserver} for execution graph export
\end{itemize}

\textbf{Source structure:}

\begin{lstlisting}[style=plain]
src/veronica_core/
  runtime_policy.py          # RuntimePolicy, PolicyContext, PolicyDecision,
                             # PolicyPipeline
  budget.py                  # BudgetEnforcer
  agent_guard.py             # AgentStepGuard
  circuit_breaker.py         # CircuitBreaker (local + distributed)
  retry.py                   # RetryContainer
  distributed.py             # RedisBudgetBackend, DistributedCircuitBreaker
  containment/
    execution_context.py     # ExecutionContext, ExecutionConfig, WrapOptions
    execution_graph.py       # ExecutionGraph, Node
    timeout_pool.py          # SharedTimeoutPool, CancellationToken
    budget_allocator.py      # BudgetAllocator (hierarchical budgets)
  middleware.py              # VeronicaASGIMiddleware, VeronicaWSGIMiddleware
  adapters/
    langchain.py             # LangChain callback adapter
    ag2.py                   # AG2 capability adapter
    langgraph.py             # LangGraph node wrapper
    mcp.py                   # MCP containment adapter (sync)
    mcp_async.py             # MCP containment adapter (async)
    crewai.py                # CrewAI adapter
    llamaindex.py            # LlamaIndex adapter
    ros2.py                  # ROS2 adapter
  partial.py                 # PartialResultBuffer
  otel.py                    # OTelExecutionGraphObserver
\end{lstlisting}

\textbf{Distribution:} \texttt{veronica-core} on PyPI (v3.4.2).

\textbf{Testing:} 4844 tests, 94\% coverage (v3.4.2). Test categories: unit
tests for each primitive, integration tests for middleware and adapters,
adversarial tests for concurrent access patterns, corrupted input handling,
and TOCTOU race conditions (\texttt{tests/adversarial/}).

\textbf{Reproducibility:} All benchmarks are deterministic (no network calls,
no random seeds beyond jitter). See \texttt{benchmarks/README.md} for
reproduction instructions.

% =====================================================================
\section{Discussion}
\label{sec:discussion}

\subsection{Design Decisions}

\textbf{Chain-level versus call-level containment.} The central design
decision in VERONICA is to enforce limits at the chain boundary rather than
per-call. This requires a persistent context object (\texttt{ExecutionContext})
that outlives individual calls, which in turn requires lifecycle management
(context managers, middleware). The alternative---per-call limits---is simpler
to implement but cannot prevent compositional failures because each individual
call is within its local budget.

\textbf{No override mechanism.} \texttt{PolicyPipeline} has no override or
exception path. Any policy denial terminates evaluation. This is intentional:
safety policies should not be overridable by lower-priority code. Applications
requiring conditional exceptions should configure narrower policies rather
than override existing ones.

\textbf{Cooperative cancellation.} G5 (wall-clock timeout) provides
cooperative rather than preemptive cancellation. This trades completeness
(in-flight operations may complete after timeout) for safety (no thread
interruption, no resource leak from killed threads). Long-running operations
can poll \texttt{CancellationToken.is\_cancelled} for finer-grained
cancellation.

\textbf{Redis for distribution.} The choice of Redis for distributed state is
pragmatic: Redis \texttt{INCRBYFLOAT} and Lua \texttt{eval()} provide the
atomic operations required for G1 and G4 in distributed mode. Alternative
backends (PostgreSQL, DynamoDB) would require explicit compare-and-swap loops;
Redis provides native atomicity at lower latency.

\subsection{Limitations}

\begin{enumerate}
  \item \textbf{Cooperative cancellation.} G5 applies only to new operations.
    In-flight operations at timeout time may complete before observing
    cancellation.

  \item \textbf{\texttt{RetryContainer} must use \texttt{ExecutionContext} for
    G3.} Standalone \texttt{RetryContainer} provides only local bounds; the
    chain-level ceiling requires an active \texttt{ExecutionContext}.

  \item \textbf{Redis availability for G1 distributed.} When Redis is
    unavailable and \texttt{fallback\_on\_error=True}, the
    \texttt{LocalBudgetBackend} fallback provides only process-local bounds,
    not cross-process bounds.

  \item \textbf{Cost accuracy.} G1 bounds accumulated reported costs, not
    actual costs. If the LLM provider reports costs inaccurately, the bound
    applies to reported values.

  \item \textbf{Agent framework integration depth.} Framework adapters wrap
    the outermost callable boundary; agent-internal retry logic (e.g.,
    LangChain's \texttt{tenacity} integration at the chain level) may not be
    captured by \texttt{RetryContainer} without explicit wrapping.

  \item \textbf{Residual cost bound.} G1 bounds the \emph{committed} cost.
    Because VERONICA uses cooperative containment (check-before-dispatch),
    calls that are already in-flight when \texttt{HALT} fires may complete and
    incur cost that is not reflected in the committed accumulator until after
    the halt. The worst-case residual cost is:
    \[
      \text{ResidualCost} \leq \text{MaxInFlight} \times \text{AvgCallCost}
    \]
    where $\text{MaxInFlight}$ is the maximum number of concurrent
    \texttt{wrap\_llm\_call()} invocations in a single
    \texttt{ExecutionContext} (bounded by the number of threads calling
    \texttt{wrap\_llm\_call()} on the same context) and $\text{AvgCallCost}$
    is the average cost of a single LLM call.

    In the common single-threaded agent pattern, $\text{MaxInFlight} = 1$,
    and the residual is at most one call's cost. In multi-threaded
    configurations with $N$ worker threads, $\text{MaxInFlight} = N$, and the
    total cost may exceed \texttt{max\_cost\_usd} by up to
    $N \times \text{AvgCallCost}$ in the worst case.

    This bound is tight: it is achievable when all $N$ threads pass the budget
    check simultaneously, each dispatches one call, and all $N$ calls succeed
    and commit cost before any observes the updated accumulator. The two-phase
    reserve/commit protocol (v2.0 roadmap) would reduce the residual to the
    cost of the single reserving call, at the expense of additional lock
    contention.

  \item \textbf{Single-language implementation.} VERONICA is implemented in
    Python. Agent systems written in other languages (TypeScript, Go, Rust)
    cannot use the library directly. The containment model is
    language-independent, but the implementation relies on Python's
    \texttt{threading.Lock} and \texttt{contextvars} for its concurrency and
    context-propagation guarantees.

  \item \textbf{GIL and concurrency model.} Python's Global Interpreter Lock
    (GIL) serializes bytecode execution, which simplifies reasoning about lock
    correctness but limits true parallelism to I/O-bound workloads. Under
    CPython with the GIL, the \texttt{threading.Lock} acquisitions in
    \texttt{BudgetEnforcer} and \texttt{AgentStepGuard} are effectively
    uncontended for CPU-bound code. On future GIL-free Python builds (PEP
    703), the lock-based protocols remain correct but may exhibit higher
    contention under true thread parallelism.

  \item \textbf{LLM cost estimation accuracy.} VERONICA's pricing table maps
    model names to per-token costs. When providers change pricing or introduce
    new models, the table may be stale. The \texttt{cost\_estimate\_hint}
    parameter in \texttt{WrapOptions} allows callers to override the estimate,
    but incorrect hints weaken the practical accuracy of G1 (the formal bound
    still holds on committed values).
\end{enumerate}

\subsection{Future Work}

The v2.0 roadmap includes: async-native budget enforcement
(\texttt{AsyncBudgetBackend}); reserve-commit-rollback budget protocol for
speculative cost accounting; \texttt{on\_halt} callback for
application-level cleanup on containment; and OpenTelemetry metric exports
for containment events. The distributed circuit breaker Lua scripts target
Redis Cluster compatibility in a future release.

% =====================================================================
\section{Conclusion}
\label{sec:conclusion}

VERONICA provides a runtime containment layer that addresses the compositional
resource amplification failures inherent in LLM agent systems. By enforcing
hard limits at the chain level rather than the call level, VERONICA provides
guarantees that per-call policies cannot: cost bounds that hold across nested
retry layers (G1), step limits that apply across tool recursion (G2), retry
budgets that span agent topologies (G3), and circuit breakers that isolate
failures across process boundaries (G4).

The four core primitives---\texttt{BudgetEnforcer}, \texttt{AgentStepGuard},
\texttt{CircuitBreaker}, and \texttt{RetryContainer}---are composable via
\texttt{PolicyPipeline} and require no modification to existing agent logic.
Empirical evaluation shows 53--90\% operation reduction versus uncontained
baselines, with full-pipeline overhead of 11.43 microseconds per call, less
than 0.002\% of typical LLM API latency. The implementation is available as
\texttt{veronica-core} on PyPI (v3.4.2, 4844 tests, 94\% coverage).

% =====================================================================
\bibliographystyle{unsrt}
\bibliography{references}

% =====================================================================
\appendix

\section{API Quick Reference}
\label{sec:appendix-api}

\begin{lstlisting}
from veronica_core.containment import ExecutionContext, ExecutionConfig, WrapOptions
from veronica_core import BudgetEnforcer, AgentStepGuard, CircuitBreaker, RetryContainer
from veronica_core.runtime_policy import PolicyPipeline, PolicyContext

# Compose a policy pipeline
pipeline = PolicyPipeline([
    BudgetEnforcer(limit_usd=10.0),
    AgentStepGuard(max_steps=25),
    CircuitBreaker(failure_threshold=5, recovery_timeout=60.0),
])

# Configure chain limits
config = ExecutionConfig(
    max_cost_usd=1.0,
    max_steps=50,
    max_retries_total=10,
    timeout_ms=30_000,
)

# Execute with containment
with ExecutionContext(config=config, pipeline=pipeline) as ctx:
    for step in agent_steps():
        decision = ctx.wrap_llm_call(
            fn=lambda: client.chat(messages=step.messages),
            options=WrapOptions(
                operation_name="agent_step",
                cost_estimate_hint=step.estimated_cost,
            ),
        )
        if decision == Decision.HALT:
            result = ctx.get_partial_result()
            break
\end{lstlisting}

\section{Benchmark Reproduction}
\label{sec:appendix-benchmarks}

\textbf{Environment.}

\begin{table}[!ht]
\centering
\begin{tabular}{@{}ll@{}}
\toprule
\textbf{Component} & \textbf{Version / Specification} \\
\midrule
Python & 3.10, 3.11, or 3.12 (tested on all three in CI) \\
OS & Linux (Ubuntu 22.04), macOS 14, Windows 11 \\
Hardware & Any x86\_64 or ARM64 (no GPU required) \\
\texttt{veronica-core} & 2.0.0 (\texttt{pip install veronica-core==2.0.0}) \\
\texttt{fakeredis} & $\geq$ 2.26 with \texttt{lupa} $\geq$ 2.6
(distributed tests only) \\
\bottomrule
\end{tabular}
\end{table}

No external services (Redis, LLM APIs) are required for benchmark
reproduction. All benchmarks use stub implementations with deterministic
behavior.

\textbf{Installation.}

\begin{lstlisting}[style=bash]
# From PyPI
pip install veronica-core==2.0.0

# From source (for development/modification)
git clone https://github.com/amabito/veronica-core
cd veronica-core
git checkout v3.4.2
pip install -e ".[dev]"
\end{lstlisting}

\textbf{Benchmark commands and expected output.}

\begin{lstlisting}[style=bash]
# Section 6.2-6.3: Four canonical failure modes
python benchmarks/bench_retry_amplification.py
# Expected: baseline_ops=27, veronica_ops=3, reduction=88.9%

python benchmarks/bench_recursive_tools.py
# Expected: baseline_ops=20, veronica_ops=5, reduction=75.0%

python benchmarks/bench_multi_agent_loop.py
# Expected: baseline_ops=60, veronica_ops=28, reduction=53.3%

python benchmarks/bench_websocket_runaway.py
# Expected: baseline_ops=100, veronica_ops=10, reduction=90.0%

# Baseline comparison (all 4 baselines x 4 scenarios)
python benchmarks/bench_baseline_comparison.py
# Expected: JSON with 16 entries, avg reduction ~78.8%

# Ablation study (component contribution)
python benchmarks/bench_ablation_study.py
# Expected: 7 configurations, monotonic reduction as components added

# Section 6.5: Scale simulation (1-1000 concurrent chains)
python benchmarks/scale_simulation.py
# Expected: ~83% reduction, ~12.6 us/chain overhead at 1000 chains

# Real incident reproductions
python benchmarks/real_incidents/run_all.py
# Expected: 5/5 PASS
\end{lstlisting}

All scripts produce JSON to stdout. Numeric results should match the values in
Section~\ref{sec:evaluation} within rounding tolerance. Timing measurements
(microseconds/call) may vary by $\pm$50\% depending on hardware; the relative
ordering and order-of-magnitude should be consistent.

\textbf{Verification against paper claims.}

\begin{table}[!ht]
\centering
\footnotesize
\caption{Paper claim verification.}
\label{tab:verification}
\begin{tabular}{@{}p{4cm}p{5cm}p{3.5cm}@{}}
\toprule
\textbf{Paper Claim} & \textbf{Verification Command} &
\textbf{Pass Criterion} \\
\midrule
88.9\% retry reduction (Section~6.3) &
\texttt{bench\_retry\_amplification.py} &
\texttt{reduction $\geq$ 0.88} \\
75.0\% recursive reduction (Section~6.3) &
\texttt{bench\_recursive\_tools.py} &
\texttt{reduction $\geq$ 0.74} \\
53.3\% multi-agent reduction (Section~6.3) &
\texttt{bench\_multi\_agent\_loop.py} &
\texttt{reduction $\geq$ 0.52} \\
90.0\% WebSocket reduction (Section~6.3) &
\texttt{bench\_websocket\_runaway.py} &
\texttt{reduction $\geq$ 0.89} \\
$<$ 12 $\mu$s/call overhead (Section~6.4) &
\texttt{bench\_baseline\_comparison.py} &
\texttt{overhead\_us $<$ 20} (hardware-dependent) \\
G1--G6 hold (Section~5) &
\texttt{pytest tests/ -q} &
4844 tests pass, 0 fail \\
\bottomrule
\end{tabular}
\end{table}

\section{Test Suite Structure}
\label{sec:appendix-tests}

\begin{lstlisting}[style=plain]
tests/
  unit/
    test_budget.py                  # BudgetEnforcer, including NaN/Inf/negative
    test_agent_guard.py             # AgentStepGuard, increment-before-check
    test_circuit_breaker.py         # State machine, HALF_OPEN slot constraint
    test_retry.py                   # RetryContainer, jitter, backoff
    test_execution_context.py       # ExecutionContext, WrapOptions, snapshots
    test_runtime_policy.py          # PolicyPipeline, PolicyContext, PolicyDecision
  integration/
    test_middleware.py              # ASGI/WSGI middleware, WebSocket
    test_adapters_langchain.py      # LangChain adapter
    test_adapters_ag2.py            # AG2 capability adapter
    test_distributed.py             # Redis backend (fakeredis + lupa)
  adversarial/
    test_concurrent_budget.py       # Race conditions on BudgetEnforcer
    test_concurrent_circuit.py      # TOCTOU on CircuitBreaker HALF_OPEN
    test_corrupted_input.py         # NaN, Inf, garbage strings in Redis state
    test_partial_failure.py         # Backend unavailable mid-operation
\end{lstlisting}

\end{document}
